\documentclass[a4paper,11pt]{article}
\usepackage[T2A]{fontenc}
\usepackage[utf8]{inputenc}
\usepackage[russian]{babel}
\usepackage{amsmath,amssymb,mathtools}
\usepackage{geometry}
\geometry{left=23mm,right=23mm,top=23mm,bottom=24mm,headheight=23pt,headsep=8pt}
\usepackage{microtype}
\usepackage{parskip}
\usepackage{setspace}
\setstretch{1.11}
\setlength{\parskip}{4pt plus 1pt minus 1pt}
\usepackage{xcolor}
\definecolor{accent}{HTML}{007F83}
\definecolor{darkaccent}{HTML}{075A5F}
\definecolor{textgray}{HTML}{414141}
\usepackage{tikz}
\usetikzlibrary{arrows.meta,positioning,shapes.geometric,calc}
\usepackage{pgfplots}
\pgfplotsset{compat=1.18}
\usepackage{tabularx,booktabs,array}
\usepackage{enumitem}
\usepackage{fancyhdr}
\usepackage{hyperref}
\hypersetup{colorlinks=true,linkcolor=darkaccent,pdfauthor={Лёвин Артём Александрович},pdftitle={Чек-лист. Задачи прикладного характера. Ярослав Гаврилов}}
\usepackage{titlesec}
\titleformat{\section}{\Large\color{darkaccent}}{\thesection.}{0.6em}{}
\titleformat{\subsection}{\large\color{darkaccent}}{}{0em}{}
\titlespacing*{\section}{0pt}{12pt}{7pt}
\titlespacing*{\subsection}{0pt}{9pt}{4pt}
\setlist[enumerate]{leftmargin=*,itemsep=8pt,topsep=5pt}
\pagestyle{fancy}
\fancyhf{}
\fancyhead[L]{\footnotesize\color{textgray}Задачи прикладного характера}
\fancyhead[R]{\footnotesize\color{textgray}Ярослав Гаврилов}
\renewcommand{\headrulewidth}{0.35pt}
\renewcommand{\headrule}{\hbox to\headwidth{\color{accent}\leaders\hrule height \headrulewidth\hfill}}
\fancyfoot[C]{\small\thepage}
\newcommand{\keyword}[1]{\textcolor{darkaccent}{#1}}
\newcommand{\idea}[1]{\par\noindent\keyword{Ключевая идея.} #1\par}
\newcommand{\example}[1]{\par\noindent\keyword{Пример.} #1\par}
\newcommand{\caution}[1]{\par\noindent\textcolor{darkaccent}{Важно.} #1\par}
\begin{document}
\thispagestyle{empty}
\begin{center}
\vspace*{17mm}
{\large\color{darkaccent}МАТЕМАТИКА\par}
\vspace{9mm}
{\LARGE\color{textgray}Чек-лист\par}
\vspace{9mm}
{\Large\color{accent}Задачи прикладного характера\\[3pt]Задание № 10 профильного ЕГЭ\par}
\vspace{13mm}
{\large Подготовка к ЕГЭ по профильной математике\par}
\vfill
{\large Ярослав Гаврилов\par}
\vspace{9mm}
{\large 5 ноября 2026 г.\par}
\vspace{9mm}
{\normalsize Лёвин Артём Александрович\\[3pt]
эксклюзивно для Ярослава Гаврилова\par}
\vspace{11mm}
\begin{tikzpicture}[x=1cm,y=1cm]
  \draw[accent,line width=.9pt] (-5,0) .. controls (-3.6,1.1) and (-2.1,-.65) .. (0,0)
  .. controls (2.1,.7) and (3.4,-1.15) .. (5,0);
  \fill[accent] (0,0) circle (1.25pt);
\end{tikzpicture}
\vspace{5mm}
\end{center}
\clearpage

\section{Главная стратегия: внимательно перевести текст в математику}
В прикладной задаче формула часто уже известна. Сложность обычно связана с тем, \emph{что именно подставлять}, как трактовать слова условия и \emph{какой из найденных корней соответствует реальной ситуации}.

\begin{enumerate}[label=\arabic*.]
\item \keyword{Определите величину в вопросе:} момент времени, длительность, расстояние, высоту, скорость или температуру.
\item \keyword{Запишите известные данные} вместе с единицами измерения. При необходимости переведите сантиметры в метры.
\item \keyword{Выберите граничное значение:} например, «высота равна 3 м», «бак пуст», «давление равно нулю».
\item \keyword{Составьте уравнение} по данной формуле и решите его.
\item \keyword{Проверьте корни по смыслу.} Время неотрицательно, скорость берут неотрицательной, а бак не может снова наполниться сам собой после опустошения.
\item \keyword{Перечитайте вопрос перед ответом.} Для продолжительности понадобится разность моментов; для момента события часто достаточно одного корня.
\end{enumerate}

\caution{Если в условии сказано «не менее», это означает знак $\ge$. Решение уравнения с равенством находит границы нужного промежутка. Затем обязательно определите, на каком промежутке выполняется неравенство.}

\section{Колодец: изменение уровня воды}
Для падения небольшого камня используется данная в задаче формула
\[
 h=5t^2,\qquad t\ge0.
\]
Здесь $h$ --- расстояние до воды, $t$ --- время падения. После дождя вода поднялась, поэтому камень летит \emph{меньше} времени.

\begin{center}
\begin{tikzpicture}[scale=.95]
  \draw[blue!65!black,line width=1pt] (0,2.15)--(0,.1)--(1.35,.1)--(1.35,2.15);
  \draw[blue!65!black,line width=1pt] (3.35,2.15)--(3.35,.1)--(4.7,.1)--(4.7,2.15);
  \draw[teal!65!black,line width=1.6pt] (0,.65)--(1.35,.65);
  \draw[teal!65!black,line width=1.6pt] (3.35,1.18)--(4.7,1.18);
  \draw[<->,>=Stealth,accent,thick] (2.4,.65)--(2.4,1.18);
  \node[right,text=darkaccent,font=\small] at (2.45,.91) {$\Delta h$};
  \node[font=\small] at (.67,-.24) {До дождя};
  \node[font=\small] at (4.02,-.24) {После дождя};
  \node[font=\small,align=center] at (.67,2.53) {$t_1$};
  \node[font=\small,align=center] at (4.02,2.53) {$t_2<t_1$};
\end{tikzpicture}
\end{center}
\[
 t_2=t_1-\Delta t,\qquad
 \boxed{\Delta h=5t_1^2-5t_2^2=5(t_1-t_2)(t_1+t_2).}
\]
\example{До дождя камень летел $0{,}6$ с; после дождя время уменьшилось на $0{,}2$ с. Тогда $t_2=0{,}4$ с и}
\[
 \Delta h=5(0{,}6^2-0{,}4^2)=5(0{,}36-0{,}16)=\boxed{1\ \text{м}}.
\]
\caution{Уменьшение времени $\Delta t$ и время после дождя $t_2$ --- разные величины.}
\clearpage

\section{Экономные вычисления: вынесение множителя и разность квадратов}
После каждой строки полезно проверять, можно ли \emph{вынести общий множитель} либо применить формулу сокращённого умножения:
\[
 a^2-b^2=(a-b)(a+b).
\]
\example{Если время падения изменилось с $1{,}2$ до $1{,}1$ с, то}
\[
\begin{aligned}
\Delta h&=5\cdot1{,}2^2-5\cdot1{,}1^2\\
&=5(1{,}2-1{,}1)(1{,}2+1{,}1)\\
&=5\cdot0{,}1\cdot2{,}3=\boxed{1{,}15\ \text{м}}.
\end{aligned}
\]
Вычисление через разность квадратов часто быстрее прямого возведения десятичных дробей в квадрат. Следите за положением запятой: $1{,}15$ и $115$ отличаются в 100 раз.

\section{Мяч выше заданной высоты: два момента и один промежуток}
Высота мяча описана формулой с урока:
\[
 h(t)=1{,}6+8t-5t^2.
\]
Требуется определить, \emph{сколько секунд} высота составляет не менее $3$ м.
Сначала найдём моменты пересечения уровня $3$ м:
\[
\begin{aligned}
1{,}6+8t-5t^2&=3,\\
5t^2-8t+1{,}4&=0,\\
25t^2-40t+7&=0 \quad (\times10,\ :2),\\
D&=(-40)^2-4\cdot25\cdot7=900,\\
t_1&=\frac{40-30}{50}=0{,}2,\qquad
 t_2=\frac{40+30}{50}=1{,}4.
\end{aligned}
\]
\begin{center}
\begin{tikzpicture}
\begin{axis}[
width=.88\linewidth,height=5.3cm,
axis lines=left,xmin=0,xmax=1.7,ymin=0,ymax=5.1,
xtick={.2,.8,1.4},xticklabels={$0{,}2$,$0{,}8$,$1{,}4$},
ytick={0,3},yticklabels={$0$,$3$},
xlabel={$t,\ \text{с}$},ylabel={$h,\ \text{м}$},
label style={font=\small},tick label style={font=\small},
every axis plot/.append style={smooth},clip=false]
\addplot[accent,line width=1.25pt,domain=0:1.65,samples=100]{1.6+8*x-5*x^2};
\addplot[black!50,densely dashed,domain=0:1.62]{3};
\addplot[darkaccent,very thick,domain=.2:1.4]{3};
\addplot[only marks,mark=*,mark size=2pt,accent] coordinates {(.2,3) (1.4,3)};
\node[font=\small,above,text=darkaccent] at (axis cs:.83,3.08) {$h\ge3$};
\end{axis}
\end{tikzpicture}
\end{center}
Между моментами $0{,}2$ и $1{,}4$ мяч находится выше уровня $3$ м, а на границах находится ровно на уровне $3$ м. Значит,
\[
\boxed{\Delta t=t_2-t_1=1{,}4-0{,}2=1{,}2\ \text{с}}.
\]
\caution{Два корня здесь обозначают \emph{начало и конец} нужного промежутка. В ответе требуется длительность, поэтому корни вычитают.}
\clearpage

\section{Опустошение бака: почему второй корень надо отбросить}
На занятии использовалась модель высоты воды
\[
H(t)=H_0-\sqrt{2gH_0}\,kt+\frac{gk^2}{2}t^2.
\]
При $H_0=20$, $g=10$, $k=\frac1{50}$ получаем
\[
H(t)=20-\frac25t+\frac{t^2}{500}.
\]
Когда останется четверть начального уровня, $H=5$:
\[
5=20-\frac25t+\frac{t^2}{500}
\quad\Longleftrightarrow\quad t^2-200t+7500=0.
\]
\[
 D=200^2-4\cdot7500=10000,\quad
 t_1=50,\quad t_2=150.
\]
Однако формула описывает реальное вытекание лишь до момента полного опустошения. Здесь
\[
H(t)=\frac{(t-100)^2}{500},\qquad H(100)=0.
\]
Вода закончится через $100$ с. Значение $t=150$ с относится к продолжению математической параболы, где высота снова растёт; для бака это лишено физического смысла. \textbf{Ответ: $50$ с.}

\idea{Сначала найдите все корни уравнения, затем проверьте физический интервал $0\le t\le t_{\text{пуст}}$. Иногда допустим только первый корень.}

\section{Вращение ведёрка: обнуление силы давления}
В верхней точке окружности сила давления воды на дно задаётся формулой
\[
P=m\left(\frac{v^2}{L}-g\right).
\]
Вода перестаёт давить на дно при $P=0$. Масса $m>0$, следовательно,
\[
 \frac{v^2}{L}-g=0,\qquad v^2=gL,\qquad
 \boxed{v=\sqrt{gL}}.
\]
\example{При $L=40$ см $=0{,}4$ м и $g=10$ м/с$^2$ получаем}
\[
v=\sqrt{10\cdot0{,}4}=\boxed{2\ \text{м/с}}.
\]
\caution{Скорость здесь --- модуль величины: отрицательный корень $-2$ не записывается как ответ. Масса $m$ на результат не влияет, поскольку $m\ne0$.}
\clearpage

\section{Температура: та же идея с двумя пересечениями}
Температура электрического прибора изменяется по закону
\[
T(t)=-t^2+14t+25,\qquad t\ge0.
\]
На протяжении какого времени температура будет \emph{не менее} $70^\circ\mathrm C$?
\[
\begin{aligned}
 -t^2+14t+25&=70,\\
 t^2-14t+45&=0,\\
 (t-5)(t-9)&=0.
\end{aligned}
\]
Парабола направлена ветвями вниз. Между $t=5$ и $t=9$ температура превышает $70^\circ\mathrm C$; на концах равна $70^\circ\mathrm C$. Поэтому
\[
\boxed{\Delta t=9-5=4\ \text{мин}}.
\]
\idea{Принцип одинаков для высоты и температуры: два пересечения горизонтального уровня, затем разность моментов. Но всегда убедитесь, что между корнями функция действительно выше уровня.}

\section{Квадратный корень из большого точного квадрата}
Когда при решении квадратного уравнения встретился крупный дискриминант, оцените \emph{между какими квадратами десятков} он лежит, затем проверьте возможные последние цифры корня.
\[
 70^2=4900,\quad 80^2=6400,\quad 90^2=8100,\quad100^2=10000.
\]
\example{Найдём $\sqrt{9409}$. Поскольку $90^2<9409<100^2$, корень лежит между $90$ и $100$. Последняя цифра числа $9409$ --- $9$, поэтому у целого корня последняя цифра может быть $3$ или $7$. Кандидаты: $93$ и $97$.}
\[
93^2=8649,\qquad97^2=9409\quad\Longrightarrow\quad
\boxed{\sqrt{9409}=97}.
\]
\caution{Такой способ рассчитан на точные квадраты целых чисел. Оценка интервала и последней цифры даёт кандидатов, а \emph{проверка возведением в квадрат} устанавливает ответ.}

\section{Краткое повторение ранее изученного}
\[
 a^m a^n=a^{m+n}\ (a>0),\qquad
 \log_a b=c\ \Longleftrightarrow\ a^c=b\ (a>0,\ a\ne1,\ b>0).
\]
В простейших тригонометрических уравнениях учитывайте все точки соответствующего уровня на единичной окружности и затем отбирайте корни на заданном промежутке.
\clearpage

\section*{Алгоритм: от формулы к ответу}
\addcontentsline{toc}{section}{Алгоритм}
\begin{center}
\vspace{8mm}
\begin{tikzpicture}[>=Stealth,thick,
  block/.style={rectangle,rounded corners=2pt,draw=accent,align=center,text width=8.9cm,minimum height=1.18cm,inner sep=5pt,font=\normalsize},
  choice/.style={diamond,aspect=3.2,draw=darkaccent,align=center,text width=5cm,inner sep=3pt,minimum height=1.25cm,font=\small},
  answer/.style={rectangle,rounded corners=2pt,draw=accent,align=center,text width=5.6cm,minimum height=1.45cm,font=\small,inner sep=4pt},
  line/.style={->,line width=.85pt,accent}]
\node[block] (a) at (0,0) {Прочитайте вопрос и определите искомую величину};
\node[block] (b) at (0,-2.0) {Переведите единицы и подставьте данные в формулу};
\node[block] (c) at (0,-4.0) {Найдите граничные моменты из уравнения};
\node[block] (d) at (0,-6.0) {Проверьте все корни по физическому смыслу};
\node[choice] (e) at (0,-8.65) {Что спрашивают?};
\node[answer] (f) at (-3.12,-12.0) {Момент события\\Выберите подходящий корень};
\node[answer] (g) at (3.12,-12.0) {Продолжительность\\Найдите границы промежутка и их разность};
\node[block] (h) at (0,-15.0) {Запишите результат и нужные единицы};
\draw[line] (a)--(b);
\draw[line] (b)--(c);
\draw[line] (c)--(d);
\draw[line] (d)--(e);
\draw[line] (e.west) -| (f.north);
\draw[line] (e.east) -| (g.north);
\node[font=\small,text=darkaccent,anchor=east] at (-3.28,-10.7) {момент};
\node[font=\small,text=darkaccent,anchor=west] at (3.28,-10.7) {длительность};
\draw[accent,line width=.85pt] (f.south) -- (-3.12,-13.65) -- (0,-13.65);
\draw[accent,line width=.85pt] (g.south) -- (3.12,-13.65) -- (0,-13.65);
\draw[line] (0,-13.65) -- (h.north);
\end{tikzpicture}
\end{center}
\clearpage

\section*{Тренировочный блок --- Путь А}
\begingroup\small
\noindent\keyword{10 заданий.} Выполняйте по порядку. В прикладных задачах внимательно прочитайте описание ситуации, определите искомую величину и учитывайте физический смысл ответа. При необходимости указывайте единицы измерения.

\subsection*{Основная тема: задачи прикладного характера}
\begin{enumerate}[label=\textbf{\arabic*.},leftmargin=1.0cm,itemsep=9pt]
\item \textbf{Повышение уровня воды в колодце.} Во время наблюдения за уровнем воды в колодце в него отпускают небольшой камень без начальной скорости. Расстояние, которое камень проходит до поверхности воды, вычисляют по формуле $h=5t^2$, где $h$~--- расстояние в метрах, а $t$~--- время падения камня в секундах. До сильного дождя время падения составляло $0{,}9$~с. После дождя поверхность воды поднялась, и камень стал достигать её на $0{,}3$~с быстрее. На сколько метров поднялся уровень воды в колодце?

\item \textbf{Извлечение квадратного корня.} При решении прикладной задачи после вычисления дискриминанта иногда требуется извлечь квадратный корень из достаточно большого целого числа. Найдите значения следующих выражений \emph{без калькулятора}. Для каждого числа сначала установите, между квадратами каких двух соседних десятков оно находится, затем проверьте найденное целое значение возведением в квадрат:
\[
\text{а)}\ \sqrt{6084};\qquad \text{б)}\ \sqrt{8649}.
\]

\item \textbf{Вращение ведёрка с водой.} Небольшое ведёрко с водой равномерно движется по вертикальной окружности на верёвке. В верхней точке траектории сила давления воды на дно ведёрка определяется формулой
\[
 P=m\left(\frac{v^2}{L}-g\right),
\]
где $m$~--- масса воды в килограммах, $v$~--- линейная скорость движения в метрах в секунду, $L$~--- радиус окружности в метрах, $g$~--- ускорение свободного падения в метрах в секунду за секунду. Известно, что длина верёвки равна $90$~см и радиус движения можно считать равным этой длине, а $g=10$~м/с$^2$. Найдите наименьшую скорость движения ведёрка, при которой в верхней точке вода перестанет давить на его дно. Ответ дайте в метрах в секунду.

\item \textbf{Полёт мяча над заданной отметкой.} Во время тренировки мяч бросили вертикально вверх с некоторой начальной высоты. Его высота над землёй меняется по закону
\[
 h(t)=1{,}6+8t-5t^2,
\]
где $h(t)$ измеряется в метрах, а $t$~--- время в секундах с момента броска. Поднимаясь, мяч впервые проходит отметку $4$~м, затем достигает наибольшей высоты и на обратном пути снова пересекает эту отметку. Определите, сколько \emph{всего секунд} от первого до второго пересечения мяч находится на высоте не менее $4$~м. Ответ дайте в секундах.

\clearpage
\item \textbf{Нагревание и остывание прибора.} После включения нагревателя температура лабораторного прибора сначала повышается, достигает максимума, а затем начинает снижаться. На рассматриваемом промежутке времени температура описывается формулой
\[
 T(t)=-t^2+12t+16,
\]
где $T(t)$~--- температура в градусах Цельсия, $t$~--- число минут, прошедших с начала наблюдения. Для проведения измерения требуется, чтобы температура прибора составляла \emph{не менее} $43^\circ\mathrm C$. Сколько минут продолжается промежуток времени, в течение которого это условие выполняется? Ответ дайте в минутах.
\item \textbf{Время полного опустошения резервуара.} Из резервуара открывают сливной кран. Высота столба воды в резервуаре в течение всего процесса опустошения изменяется по формуле
\[
 H(t)=9-\frac35t+\frac{t^2}{100},
\]
где $H(t)$ выражена в метрах, $t$~--- время в секундах с момента открытия крана. Формула применима, пока в резервуаре есть вода; после его опустошения уровень остаётся нулевым. Определите, через сколько секунд после открытия крана резервуар станет полностью пустым. Ответ дайте в секундах.

\item \textbf{Четверть первоначального объёма воды.} В резервуаре с постоянной площадью горизонтального сечения открыли сливной кран. По мере вытекания воды её уровень изменяется по формуле
\[
 H(t)=18-\frac35t+\frac{t^2}{200},
\]
где $H(t)$~--- высота воды в метрах, $t$~--- число секунд с момента открытия крана. В начальный момент резервуар заполнен до уровня $18$~м. Так как площадь сечения постоянна, объём оставшейся воды пропорционален её высоте. Через сколько секунд после открытия крана объём воды \emph{впервые} составит четверть первоначального? Учитывайте, что после полного опустошения формула перестаёт описывать реальную ситуацию. Ответ дайте в секундах.

\item \textbf{Максимальная высота и граничное значение.} Мяч бросили вертикально вверх с площадки, расположенной на высоте $2$~м над землёй. На рассматриваемом промежутке времени его высота определяется формулой
\[
 h(t)=2+6t-2t^2,
\]
где $h(t)$ измеряется в метрах, а $t$~--- время в секундах с момента броска. Требуется установить, в какой момент мяч достигает наибольшей высоты. Кроме того, необходимо выяснить, сколько секунд мяч находится на высоте \emph{не менее $6{,}5$~м}. Для ответа на второй вопрос учитывайте, может ли указанная высота достигаться в течение целого промежутка времени или только в отдельный момент.
\end{enumerate}

\subsection*{Повторение: степени, логарифмы, тригонометрия}
\begin{enumerate}[label=\textbf{\arabic*.},leftmargin=1.0cm,itemsep=9pt,start=9]
\item Вычислите:
\[
\frac{2^5\cdot 2^{-2}}{4}+\log_3 81.
\]

\item Решите уравнение $\cos x=\frac12$ на промежутке $0\le x<2\pi$.
\end{enumerate}

\vspace{7mm}
\noindent\keyword{Перед сдачей:} сопоставьте единицы измерения; проверьте знаки в квадратном уравнении, все полученные корни и смысл вопроса. Для слов «не менее» определите \emph{весь} нужный промежуток.
\par\endgroup
\clearpage

\section*{Ответы к тренировочному блоку}
\vspace{7mm}
\renewcommand{\arraystretch}{1.35}
\begin{center}
\begin{tabularx}{\linewidth}{@{}>{\bfseries}p{13mm}X@{}}
\toprule
№ & Ответ \\
\midrule
1 & $2{,}25$ м \\
2 & а) $78$; б) $93$ \\
3 & $3$ м/с \\
4 & $0{,}8$ с \\
5 & $6$ мин \\
6 & $30$ с \\
7 & $30$ с \\
8 & Максимум при $t=1{,}5$ с; продолжительность $0$ с \\
9 & $6$ \\
10 & $\dfrac{\pi}{3}$; $\dfrac{5\pi}{3}$ \\
\bottomrule
\end{tabularx}
\end{center}

\vspace{9mm}
\section*{Итоговая памятка}
\renewcommand{\arraystretch}{1.28}
\begin{tabularx}{\linewidth}{@{}p{.36\linewidth}X@{}}
\toprule
Сигнал в условии & Действие \\
\midrule
«Время изменилось на...» & Вычислите новое время, затем разность высот.\\
«Не менее...» & Найдите граничные моменты и проверьте промежуток между ними.\\
«Через сколько...» & Выберите физически допустимый момент события.\\
«Бак пуст» & Приравняйте высоту к нулю; после опустошения формула может не работать.\\
Большой точный квадрат & Оцените десятки, отберите кандидатов и проверьте квадрат.\\
Масса в ненулевом множителе & Приравняйте к нулю нужный второй множитель.\\
\bottomrule
\end{tabularx}
\vfill
\begin{center}\footnotesize\color{textgray}Лёвин Артём Александрович эксклюзивно для Ярослава Гаврилова\end{center}
\end{document}
